\honest{How Much Evidence Does It Take?}{Eliezer Yudkowsky}{2007}

How much evidence does it take to believe something? I start with a lottery: 131,115,985 possible tickets, one winner, all equally likely.

A black box always beeps for the winning ticket, and beeps a quarter of the time for a losing one. Tried on every ticket, it beeps for about 33 million losers. Two such boxes, independent of each other, cut that to about 8 million. Evidence 131 million times more likely for the winner than for a loser leaves you with the winner and, on average, one loser, so an even chance. All of this is correct and clearly worked out.

Then I measure evidence in bits. I define a bit by the information in an outcome: learning that something with probability 1/8 happened gives you three bits. Then I add up the bits from the boxes a different way, two per box, because each box makes a beep four times more likely for the winner. The two measures agree in the lottery, where nearly every ticket loses. Elsewhere they do not, and I do not say so.

Twenty-eight bits, or 14 boxes, bring you to odds of 2 to 1. Thirty-four bits make you 99\% sure. Along the way I call a probability of fifty percent ``Fifty percent odds,'' and call odds of 128 to 1 ``the chance of winning,'' in a post that spends a sentence explaining the difference between odds and probability.

Then the general rule: the more possibilities there are, the less likely the hypothesis ``seems a priori,'' and the surer you want to be, the more evidence you need. In the lottery the starting probability was given exactly. In any real question, setting it is the hard part, and ``seems'' is all I say about it.

``You cannot form accurate beliefs based on inadequate evidence.'' A lucky guess is accurate. What thin evidence cannot give you is a right to be confident.

Believing without enough evidence, I say, is like driving a car without fuel. A car without fuel does not move at all. A belief without enough evidence can still happen to be right.
